\section{Background}
\label{sec:background}



\section{Background}
\label{sec:background}

\subsection{(Extended) test martingales and (extended) Ville inequalities}
\label{subsec:ville}

Let
\[
  \bigl(\Omega,\mathcal A,(\mathcal F_n)_{n\geq 1}\bigr)
\]
be a filtered measurable space, and let \(\mathcal P\) be a family of
probability measures on \((\Omega,\mathcal A)\), representing a
composite null hypothesis. We allow a process to start at an arbitrary
deterministic time \(m\geq 1\); the usual case is \(m=1\).

\begin{definition}[\(p\)-process]
An adapted process \((P_n)_{n\geq m}\), taking values in \([0,1]\), is
called a \(p\)-process for \(\mathcal P\) if, for every
\(\mathbb P\in\mathcal P\) and every \(\alpha\in[0,1]\),
\[
  \mathbb P\left(\inf_{n\geq m}P_n\leq\alpha\right)
  \leq\alpha.
\]
\end{definition}

\begin{definition}[Test martingale]
An adapted process \((X_n)_{n\geq m}\) is a test martingale for
\(\mathcal P\), with respect to \((\mathcal F_n)_{n\geq m}\), if, for
every \(\mathbb P\in\mathcal P\), the random variables \(X_n\) are
nonnegative and \(\mathbb P\)-integrable,
\[
  \mathbb E_{\mathbb P}[X_m]=1,
\]
and
\[
  \mathbb E_{\mathbb P}
  \left[X_{n+1}\mid\mathcal F_n\right]
  =
  X_n
  \qquad
  \mathbb P\text{-a.s. for every }n\geq m.
\]
\end{definition}

The integrability requirement can be relaxed by considering extended
nonnegative martingales \citep{Wang2025a}.

\begin{definition}[Extended nonnegative martingale]
An adapted process \((X_n)_{n\geq m}\), taking values in
\([0,\infty]\), is an extended nonnegative martingale for
\(\mathcal P\), with respect to \((\mathcal F_n)_{n\geq m}\), if, for
every \(\mathbb P\in\mathcal P\),
\[
  \mathbb E_{\mathbb P}
  \left[X_{n+1}\mid\mathcal F_n\right]
  =
  X_n
  \qquad
  \mathbb P\text{-a.s. for every }n\geq m,
\]
where the conditional expectation is understood in the extended
nonnegative sense.
\end{definition}

Every test martingale is an extended nonnegative martingale.
Conversely, under a given \(\mathbb P\in\mathcal P\), an extended
nonnegative martingale is a classical nonnegative martingale if
\(X_m\) is \(\mathbb P\)-integrable; it is a test martingale if, in
addition, \(\mathbb E_{\mathbb P}[X_m]=1\). The extended nonnegative
martingales used below are finite \(\mathbb P\)-almost surely from
their stated starting times, but need not be integrable.

The following result is a direct consequence of Ville's inequality.

\begin{proposition}[Ville calibration of a test martingale]
\label{prop:ville-calibration}
Let \((X_n)_{n\geq m}\) be a test martingale for \(\mathcal P\).
Define
\[
  P_n
  =
  \min\left\{1,\frac{1}{X_n}\right\},
  \qquad
  \widetilde P_n
  =
  \min\left\{
    1,
    \frac{1}{\max_{m\leq k\leq n}X_k}
  \right\},
  \qquad n\geq m,
\]
where \(1/0:=\infty\). Then
\((P_n)_{n\geq m}\) and
\((\widetilde P_n)_{n\geq m}\) are \(p\)-processes for
\(\mathcal P\).
\end{proposition}

We call a \(p\)-process obtained in this way a
Ville-calibrated \(p\)-process.

The corresponding result for extended nonnegative martingales follows
from the extended Ville inequality for extended nonnegative
supermartingales \citep{Wang2025a}.

\begin{proposition}[Extended Ville calibration]
\label{prop:extended-ville-calibration}
Let \((X_n)_{n\geq m}\) be an extended nonnegative martingale for
\(\mathcal P\). For \(x>0\), define
\[
  g_m(x)
  =
  \sup_{\mathbb P\in\mathcal P}
  \mathbb E_{\mathbb P}
  \left[
    \min\left\{1,\frac{X_m}{x}\right\}
  \right],
\]
and set \(g_m(0)=1\). Then
\[
  \bigl(g_m(X_n)\bigr)_{n\geq m}
\]
and
\[
  \left(
    g_m\left(\max_{m\leq k\leq n}X_k\right)
  \right)_{n\geq m}
\]
are \(p\)-processes for \(\mathcal P\).
\end{proposition}

We call a \(p\)-process obtained in this way an extended
Ville-calibrated \(p\)-process. This calibration applies to extended
nonnegative martingales that need not be integrable. If
\((X_n)_{n\geq m}\) is a test martingale, then
\[
  g_m(x)
  \leq
  \min\left\{1,\frac{1}{x}\right\}.
\]
Consequently, the extended Ville-calibrated \(p\)-process is
pointwise no larger than, and hence dominates, the corresponding
Ville-calibrated \(p\)-process.


\subsection{Martingale construction for the anytime-valid \(t\)-test}
\label{subsec:martingale-t-test}

Suppose that
\[
  Y_1,Y_2,\ldots
  \stackrel{\mathrm{iid}}{\sim}
  \mathcal N(\mu,\sigma^2),
  \qquad \sigma>0,
\]
and define the standardized effect size
\[
  \theta^\star=\frac{\mu}{\sigma}.
\]
We consider the composite null hypothesis
\[
  H_0:\theta^\star=0,
  \qquad
  \mathcal P_0
  =
  \left\{
    \mathbb P_{0,\sigma}:\sigma>0
  \right\}.
\]

Following the constructions of
\citet{Lai1976,PerezOrtiz2024,Wang2025}, define the restricted
scale-invariant filtration
\[
  \mathcal F_n
  =
  \sigma\left(
    \frac{Y_1}{|Y_1|},
    \frac{Y_2}{|Y_1|},
    \ldots,
    \frac{Y_n}{|Y_1|}
  \right).
\]
The quantities in this display can be defined arbitrarily on the
event \(\{Y_1=0\}\), which has probability zero under every model
considered here.

Let
\[
  Z_i=\frac{Y_i}{|Y_1|},
  \qquad i\geq 1.
\]
For \(\theta\in\mathbb R\), let \(\eta_{n,\theta}\) denote the density
of \((Z_1,\ldots,Z_n)\) when \(\theta^\star=\theta\), with respect to
counting measure on \(\{-1,1\}\) in the first coordinate and
Lebesgue measure on \(\mathbb R^{n-1}\) in the remaining coordinates.
This distribution depends on \(\theta\), but not on the unknown scale
parameter \(\sigma\).

The scale-invariant likelihood ratio
\begin{equation}
\label{eq:scale-inv-LR}
  L_{n,\theta}(Y_{1:n})
  =
  \frac{
    \eta_{n,\theta}
    \left(
      Y_1/|Y_1|,
      \ldots,
      Y_n/|Y_1|
    \right)
  }{
    \eta_{n,0}
    \left(
      Y_1/|Y_1|,
      \ldots,
      Y_n/|Y_1|
    \right)
  }
\end{equation}
is therefore a \((\mathcal F_n)\)-test martingale under every
\(\mathbb P_{0,\sigma}\in\mathcal P_0\).

A direct computation \cite[Sec.~4.4]{Wang2025} gives
\begin{equation}
\label{eq:hthetadef}
  L_{n,\theta}(Y_{1:n})
  =
  e^{-n\theta^2/2}
  A_n\bigl(\theta u_n(Y_{1:n})\bigr),
\end{equation}
where
\begin{align}
  A_n(z)
  &=
  \frac{1}{\Gamma(n/2)}
  \int_0^\infty
  y^{n/2-1}
  \exp\left(-y+z\sqrt{2y}\right)
  \dd y
  \notag\\
  &=
  \mathbb E\left[e^{zR_n}\right],
  \qquad
  R_n\sim\chi_n,
  \label{eq:Andef}
\end{align}
and
\begin{equation}
\label{eq:undef}
  u_n(Y_{1:n})
  =
  \frac{\sum_{i=1}^nY_i}
       {\sqrt{\sum_{i=1}^nY_i^2}}.
\end{equation}
For \(n\geq2\), this statistic can equivalently be written as
\begin{equation}
\label{eq:undef-tstat}
  u_n(Y_{1:n})
  =
  \frac{n\widehat\theta_n}
       {\sqrt{(n-1)+n\widehat\theta_n^2}}
  =
  \frac{\sqrt n\,T_n}
       {\sqrt{(n-1)+T_n^2}},
\end{equation}
where
\[
  \overline Y_n
  =
  \frac1n\sum_{i=1}^nY_i,
  \qquad
  S_n^2
  =
  \frac{1}{n-1}
  \sum_{i=1}^n
  (Y_i-\overline Y_n)^2,
\]
\[
  \widehat\theta_n
  =
  \frac{\overline Y_n}{S_n},
  \qquad
  T_n
  =
  \sqrt n\,\widehat\theta_n.
\]
Thus, \(\widehat\theta_n\) is the sample standardized effect size, or
sample Sharpe ratio, and \(T_n\) is the classical one-sample
\(t\)-statistic.


\subsection{Mixture construction}
\label{subsec:mixture-construction}

Let \(\Pi\) be a \(\sigma\)-finite measure on \(\mathbb R\),
representing a mixing measure over the effect size. When \(\Pi\) is a
probability distribution, it can also be interpreted as a prior
distribution. Define
\begin{align}
  M_n^\Pi(Y_{1:n})
  &=
  \int_{\mathbb R}
  L_{n,\theta}(Y_{1:n})
  \,\Pi(\dd\theta)
  \notag\\
  &=
  \int_{\mathbb R}
  e^{-n\theta^2/2}
  A_n\bigl(\theta u_n(Y_{1:n})\bigr)
  \,\Pi(\dd\theta).
  \label{eq:Mn-def}
\end{align}
As the integrand is nonnegative, the integral is well-defined as an
element of \([0,\infty]\).

If \(\Pi\) is a proper probability distribution, then
\[
  \bigl(M_n^\Pi(Y_{1:n})\bigr)_{n\geq1}
\]
is a nonnegative \((\mathcal F_n)\)-test martingale under
\(\mathcal P_0\). Consequently,
\begin{equation}
\label{eq:proper-mixture-p-process}
  P_n^{V,\Pi}
  =
  \min\left\{
    1,
    \frac{1}{M_n^\Pi}
  \right\}
\end{equation}
and
\begin{equation}
\label{eq:proper-mixture-running-p-process}
  \widetilde P_n^{V,\Pi}
  =
  \min\left\{
    1,
    \frac{1}{\max_{1\leq k\leq n}M_k^\Pi}
  \right\}
\end{equation}
are Ville-calibrated \(p\)-processes.

For a general \(\sigma\)-finite mixing measure \(\Pi\), the process
\((M_n^\Pi)\) is an extended nonnegative martingale
\citep{Wang2025a}. Let \(m\) be a deterministic starting time from
which \(M_n^\Pi\) is finite almost surely, and define, for \(c>0\),
\begin{equation}
\label{eq:general-extended-calibration}
  g_{\Pi,m}(c)
  =
  \sup_{\sigma>0}
  \mathbb E_{0,\sigma}
  \left[
    \min\left\{
      1,
      \frac{M_m^\Pi}{c}
    \right\}
  \right],
\end{equation}
with \(g_{\Pi,m}(0)=1\). Because \(M_m^\Pi\) is scale invariant, its
null distribution does not depend on \(\sigma\), so the supremum in
\eqref{eq:general-extended-calibration} may equivalently be replaced
by an expectation under any fixed \(\mathbb P_{0,\sigma}\). The
processes
\[
  P_n^{E,\Pi}
  =
  g_{\Pi,m}(M_n^\Pi)
\]
and
\[
  \widetilde P_n^{E,\Pi}
  =
  g_{\Pi,m}
  \left(
    \max_{m\leq k\leq n}M_k^\Pi
  \right),
  \qquad n\geq m,
\]
are extended Ville-calibrated \(p\)-processes.


\subsection{Connection to Bayes factors for Bayesian \(t\)-tests}
\label{subsec:bayesian-t-test}

The mixture process \(M_n^\Pi\) has a direct Bayesian \(t\)-test
interpretation. Parameterize the Gaussian model as
\[
  Y_i\mid(\theta,\sigma)
  \stackrel{\mathrm{iid}}{\sim}
  \mathcal N(\sigma\theta,\sigma^2),
  \qquad
  \theta=\frac{\mu}{\sigma}.
\]
Under both hypotheses, assign the nuisance scale the Jeffreys prior
\[
  \pi(\dd\sigma)\propto\frac{\dd\sigma}{\sigma}.
\]
Under \(H_0\), set \(\theta=0\), while under \(H_1\), assign the
standardized effect size the proper prior \(\theta\sim\Pi\),
independently of \(\sigma\).

The corresponding Bayes factor in favor of the alternative is
\begin{align}
  \mathrm{BF}_{10,n}^{\Pi}(y_{1:n})
  &=
  \frac{
    \displaystyle
    \int_{\mathbb R}\int_0^\infty
    p_{\theta,\sigma}^{(n)}(y_{1:n})
    \frac{\dd\sigma}{\sigma}\,
    \Pi(\dd\theta)
  }{
    \displaystyle
    \int_0^\infty
    p_{0,\sigma}^{(n)}(y_{1:n})
    \frac{\dd\sigma}{\sigma}
  }
  \notag\\
  &=
  \int_{\mathbb R}
  L_{n,\theta}(y_{1:n})\,\Pi(\dd\theta)
  =
  M_n^\Pi(y_{1:n}).
\end{align}
The arbitrary multiplicative constant in the improper scale prior
cancels because the same prior is used under both hypotheses.
Consequently,
\[
  M_n^\Pi
  =
  \mathrm{BF}_{10,n}^{\Pi}
  =
  \left(\mathrm{BF}_{01,n}^{\Pi}\right)^{-1}.
\]
Thus, when the Bayes factor is defined in favor of the null, the test
martingale is its inverse. 

Standard choices of \(\Pi\) include Gaussian effect-size priors
\citep{Gonen2005}, the centered Cauchy prior used in the
Jeffreys--Zellner--Siow Bayesian \(t\)-test \citep{Rouder2009}, and,
more generally, location-scale Student \(t\) priors
\citep{Gronau2020}. One-sided tests can be obtained by truncating and
renormalizing these priors to either positive or negative effect
sizes.

This Bayes-factor interpretation requires \(\Pi\) to be proper. For a
merely \(\sigma\)-finite mixing measure, such as Lebesgue measure,
\(M_n^\Pi\) remains an extended nonnegative martingale, but it is not
a well-defined Bayes factor because the arbitrary normalization of
\(\Pi\) does not cancel.


\subsection{Lebesgue mixing measure}
\label{subsec:improper-prior}

The improper Lebesgue mixing measure
\[
  \Pi_{\mathrm{Leb}}(\dd\theta)
  =
  \dd\theta
\]
is an important noninformative benchmark
\citep{Lai1976,Wang2025}. In this case,
\[
  M_1^{\mathrm{Leb}}=\infty
  \qquad
  \mathbb P_{0,\sigma}\text{-a.s.},
\]
whereas \(M_n^{\mathrm{Leb}}\) is finite almost surely for every
\(n\geq2\). The process
\[
  \bigl(M_n^{\mathrm{Leb}}\bigr)_{n\geq2}
\]
is therefore a finite-valued extended nonnegative martingale, but is
not integrable and hence is not a test martingale. Standard Ville
calibration cannot be applied, but the extended Ville calibration
can be started at \(m=2\).

The martingale and its calibration function admit explicit
expressions \citep{Wang2025}.

\begin{proposition}
\label{thm:Mn-Leb}
For every \(n\geq2\),
\begin{equation}
\label{eq:Mn-Leb}
  M_n^{\mathrm{Leb}}(Y_{1:n})
  =
  \sqrt{\frac{2\pi}{n}}
  \left(
    1+
    \frac{n\widehat\theta_n^2}{n-1}
  \right)^{n/2}.
\end{equation}
Define \(g_{\mathrm{Leb}}(0)=1\) and, for \(c>0\),
\begin{equation}
\label{eq:g-Leb}
  g_{\mathrm{Leb}}(c)
  =
  \begin{cases}
    1,
    & 0<c\leq\sqrt\pi,
    \\[8pt]
    \displaystyle
    \frac{2}{\pi}
    \left[
      \arcsin\left(
        \sqrt{\frac{\sqrt\pi}{c}}
      \right)
      +
      \sqrt{
        \frac{\sqrt\pi}{c}
        \left(
          1-\frac{\sqrt\pi}{c}
        \right)
      }
    \right],
    & c>\sqrt\pi.
  \end{cases}
\end{equation}
Then
\[
  P_n^{E,\mathrm{Leb}}
  =
  g_{\mathrm{Leb}}
  \left(
    M_n^{\mathrm{Leb}}
  \right)
\]
and
\[
  \widetilde P_n^{E,\mathrm{Leb}}
  =
  g_{\mathrm{Leb}}
  \left(
    \max_{2\leq k\leq n}
    M_k^{\mathrm{Leb}}
  \right),
  \qquad n\geq2,
\]
are \(p\)-processes for the composite null family
\(\mathcal P_0\).
\end{proposition}